\subsection{S1 Probe Signal and Initial Escalation (Ablation A1)}
\label{sec:s1-probe-ablation}
% Role: component and end-to-end evidence for the S1 probe signal on UDIS-D.
\begin{table}[t]
\centering
\caption{Ablation A1 on the UDIS-D test set ($1{,}105$ pairs, $512\times 512$ RGB)
on one NVIDIA Tesla T4, with 2 warmups and 10 measured repeats per pair.
Times are in milliseconds, summarized across per-pair medians.
$t_{\text{probe}}$ is the GPU time for $S_{\text{probe}}=0.25$ downsampling and
gradient reduction across both images; $t_{\text{ctrl}}$ includes stream
synchronization and host decision time.}
\label{tab:s1-a1}
\small
% Generated by tools/s1/a1_report.py; do not edit by hand.
\begin{tabular}{lrrrrrrr}
\hline
Variant & Pairs & Retry (\%) & Att./pair & $t_{\text{probe}}$ & $t_{\text{ctrl}}$ & $t_{\text{e2e}}$ med. & $t_{\text{e2e}}$ p95 \\
\hline
Adaptive with Probe & 1105 & 18.5 & 1.363 & 0.054 & 0.064 & 4.759 & 8.598 \\
Blind Escalation ($S_0=0.25$) & 1105 & 37.6 & 1.738 & 0.000 & 0.001 & 2.308 & 6.498 \\
Blind Escalation ($S_0=0.50$) & 1105 & 16.8 & 1.362 & 0.000 & 0.001 & 4.570 & 8.525 \\
Conservative High ($S=1.0$) & 1105 & 0.0 & 1.000 & 0.000 & 0.000 & 9.020 & 17.644 \\
\hline
\end{tabular}

\end{table}
Table~\ref{tab:s1-a1} evaluates the S1 gradient probe ($S_{\text{probe}}=0.25$)
against two probe-free escalation baselines (\texttt{--no-probe}) and a fixed
high-budget baseline ($S=1.0$, $K=2000$, $R=1024$) across the $1{,}105$ pairs of
the UDIS-D test archive under the 5-rung \texttt{legacy-v1} ladder:
Rung~0 ($0.25, 500, 128$), Rung~1 ($0.50, 1000, 128$), Rung~2 ($0.50, 1000, 256$),
Rung~3 ($0.75, 2000, 512$), and Rung~4 ($1.00, 2000, 1024$).

Separating the probe with dedicated CUDA events shows that the GPU probe cost is
small. Across the $1{,}105$ pairs, $4\times 4$ area downsampling of both input
images to $128\times 128$ takes a median of $0.038$\,ms (p95 $0.044$\,ms), and
the spatial gradient reduction takes a median of $0.015$\,ms (p95 $0.016$\,ms).
Together, the GPU probe requires $0.053$\,ms ($1.1\%$ of the $4.716$\,ms median
end-to-end time), while stream synchronization and host selection add another
$0.009$\,ms ($t_{\text{ctrl}}$ median $0.062$\,ms). When the controller selects
Rung~0 ($S_0=0.25$), the first attempt reuses the probe grayscale buffers without
rerunning S1.

Compared with Conservative High ($S=1.0$), starting at $S=0.50$ via the probe
cuts median end-to-end latency from $8.931$ to $4.716$\,ms (a paired median
reduction of $4.830$\,ms, or $2.12\times$) and reduces p95 latency from $17.497$
to $8.512$\,ms, at the cost of rejecting $161$ pairs that fail the online
self-checks at the top rung. Compared with Blind Escalation starting at Rung~0
($S_0=0.25$), the probe halves the retry rate from $37.6\%$ ($416$ pairs) to
$18.5\%$ ($204$ pairs) and lowers the mean attempt count from $1.738$ to
$1.363$. On the $456$ low-gradient pairs ($g < 0.04$), blind $S_0=0.25$ retries
on $63.6\%$ of pairs ($2.333$ attempts per pair), whereas the probe skips to
Rung~2 and retries on $37.1\%$ ($1.697$ attempts).

Yet Blind Escalation at $S_0=0.25$ achieves a lower overall median latency
($2.286$\,ms vs.\ $4.716$\,ms), and Blind Escalation at $S_0=0.50$ is slightly
faster ($4.530$\,ms, $16.8\%$ retries). The trace exposes two policy causes.
First, on UDIS-D ($512\times 512$), the normalized probe gradient
$g=\min(G_A,G_B)$ spans $[0.004, 0.087]$ with median $0.044$ and 90th percentile
$0.062$. Because the uncalibrated \texttt{legacy-v1} rule requires $g\ge 0.08$
for Rung~0, the probe routes only $3$ of $1{,}105$ pairs ($0.3\%$) to $S_0=0.25$
and sends $99.7\%$ to $S_0=0.50$, discarding the $0.25$-scale probe images and
quadrupling the pixel area for S1--S6. By contrast, when $S_0=0.25$ is tried
unconditionally, $689$ pairs ($62.4\%$) pass Rung~0's loose self-check
($15$ inliers, $0.15$ ratio, $4.0$\,px residual) on the first attempt in a
median of $2.183$\,ms. Second, for the $456$ pairs with $g<0.04$, jumping from
Rung~1 to Rung~2 raises $R$ from $128$ to $256$ while tightening the self-check
from $(20\text{ inliers}, 0.15\text{ ratio}, 3.5\,\text{px})$ to
$(25, 0.20, 3.0\,\text{px})$, causing $18$ additional pairs to escalate to
$S\ge 0.75$. Thus, while the GPU probe itself costs only $0.053$\,ms, converting
that signal into an end-to-end speedup requires calibrating the gradient cutoffs
to the workload distribution so that passing $S=0.25$ pairs are not prematurely
promoted to $S=0.50$.
